\section{Materials and Methods}\label{sec:methods}
\subsection{Overview of the Two Estimators}
Figure~\ref{fig:workflow} places the original method and PPK-FIP side by side. In the original two-step method (panel a) the profiles are first plotted against the degree of reaction $c$ and fitted with the irreversible fingerprint model, which yields $S_M$ and $S_D$; these are then held fixed while the dynamic model is fitted against time in one local search, and the ``simplified'' model is chosen by comparing in-sample $F_\mathrm{obj}$ and BIC. In PPK-FIP (panel b) the composition path in $c$ is generalised to include pseudo-processivity, an exact time map adds the lumped constants and within-batch deactivation, and a single maximum-a-posteriori (MAP) fit estimates all parameters at once from global multistarts; the fingerprint estimate is not discarded but enters as the centre of a calibrated prior (dashed arrow), while shrinkage priors keep the new mechanisms at zero unless the data demand them. Two model forms are then averaged with BIC weights, and the prior hyper-parameters are chosen by meta-validation across the other case studies. The three blocks of panel (b) answer, in turn, the three open points of the Introduction: the generalised path addresses pseudo-processivity (i), the time map addresses within-batch deactivation (ii), and the prior and the out-of-sample protocols in the bottom box address the fix-then-fit constraint and the in-sample-only assessment (iii). Both estimators feed the same evaluation battery, which implements Steps~1--6 of Table~\ref{tab:roadmap}.
\begin{figure}[!htbp]\centering\incfig{fig01_workflow}\caption{Original two-step method (a) and the proposed PPK-FIP (b). The dashed arrow indicates that the fingerprint estimate, a fixed value in (a), becomes the centre of a calibrated prior in (b). Both estimators are evaluated with the same protocols (bottom).}\label{fig:workflow}\end{figure}

\subsection{Data and Audit Protocol}
All reaction profiles were taken programmatically (no manual typing) from Tables S1--S6 and S17--S21 of the Supplementary Material of \cite{nishimura2026}: Åkerman et al. at 70 and 100\,°C (oleic acid, Novozym 435, 3:1) \cite{akerman2011}; the first experiment of Bornadel et al. (same system) \cite{bornadel2013}; Tao et al. at water activities 0.25, 0.35 and 0.45 (caprylic acid, lipase of \emph{Candida} sp.\ 99--125, 3.2:1) \cite{tao2015}; and the first three reuse batches of Mao et al. (pelargonic acid, same lipase) \cite{mao2024}. Species are molar percentages relative to the initial TMP, and $c=M+2D+3T$ is the percentage degree of reaction. In total there are \NRows{} rows in 11 profiles, of which the \NFitProfiles{} profiles used in the paper contain \NFitPoints{} fitted observations (Table~\ref{tab:data}).
\inctab{data}

Figure~\ref{fig:data} shows the nine fitted profiles against time (left panel of each pair) and against $c$ (right panel). Three features of the data recur throughout the paper. First, the sampling density differs by a factor of four between studies, from 7--9 sampling times in the Åkerman profiles to 30 in Bornadel; the sparse Åkerman profiles are the ones in which prior information matters most (Step~4). Second, against $c$ the three Tao water activities nearly coincide, and so do the three Mao batches, although their time-scales differ: this is the empirical basis of the fingerprint idea, that the composition path is a property of the enzyme--substrate pair while the time-scale depends on the conditions, and it is the property that Step~5 tests by transfer. Third, the Mao data lack Z, and in batch~1 the triester reaches a plateau near 90\,\% instead of 100\,\%, a feature that mass-conserving models treat, by construction, as a property of the digitised data rather than of the kinetics (Steps~2 and 6).
\begin{figure}[!htbp]\centering\incfig{fig02_data_overview}\caption{The nine reaction profiles (Supplementary Tables S1--S6, S17--S19 of \cite{nishimura2026}) against time (left of each pair) and against the degree of reaction $c$ (right). Z, M, D, T: TMP and its mono-, di- and triester (molar \% of initial TMP).}\label{fig:data}\end{figure}

Because the paper does not print sample sizes, the public data were checked against every printed quantity that depends on them: the BIC of all twelve dynamic fits, which can only be reproduced with the exact number of data points $n_T$, and the sum of squares of the four fingerprint fits, which requires every individual value to match. The printed dynamic parameters were then evaluated in our implementation (``at printed $\theta$''), the same models were re-fitted from 30 multistarts (``global''), 200 random starts of the full model were used to map its local optima, and the profile objective along $F_A$ (all other parameters re-optimised) was computed to test identifiability \cite{raue2009}.

\subsection{Reference Model: the Original Two-Step Method}
For irreversible esterification with water removed, the fingerprinting model of \cite{nishimura2026} is
\begin{equation}
\begin{aligned}
&\frac{dZ}{dc}=\frac{-Z}{S},\qquad \frac{dM}{dc}=\frac{Z-S_MM}{S},\qquad \frac{dD}{dc}=\frac{S_MM-S_DD}{S},\qquad\frac{dT}{dc}=\frac{S_DD}{S},\\
&S=Z+S_MM+S_DD,
\end{aligned}\label{eq:fp}
\end{equation}
and the dynamic model shares the Ping Pong denominator
\begin{equation}
\begin{aligned}
&r=\frac{k_A[E]_T\,A}{A(F_A+F_{AZ}Z+F_{AM}M+F_{AD}D)+Z+S_MM+S_DD},\\
&\frac{dZ}{dt}=-rZ,\quad \frac{dM}{dt}=r(Z-S_MM),\quad \frac{dD}{dt}=r(S_MM-S_DD),\quad \frac{dT}{dt}=rS_DD,
\end{aligned}\label{eq:dyn}
\end{equation}
with $A=A_0-c$ the fatty acid. Step 1 estimates $S_M,S_D$ by least squares on profiles plotted against $c$; step 2 fixes them and estimates $k_E=k_A[E]_T$ and the lumped constants against time; the ``simplified'' model sets the constants that appear unimportant to zero. Because $k_A$ and $[E]_T$ only appear as a product, we estimate $k_E$ and recover the enzyme loading as $k_E/k_A$; the loadings implied by the printed tables are round numbers (Section~\ref{sec:repro}), which confirms that the parameterisation is identical. The objective is $F_\mathrm{obj}=\sum(y_\mathrm{pred}-y_\mathrm{obs})^2$ over Z, M, D, T (Z omitted for Mao) and $\mathrm{BIC}=n_T\ln(F_\mathrm{obj}/n_T)+k\ln n_T$ \cite{schwarz1978}. Following the paper, $k$ counts only the parameters of step 2 (``paper convention''); we additionally report $\mathrm{BIC}_\mathrm{tot}$, in which the two selectivities estimated from the same data are also counted. Throughout, ``ORIG'' denotes the paper's simplified model and ``ORIG-full'' its full model, both re-fitted by us to their global optimum.

\subsection{Generalised Mechanism: Pseudo-Processivity and Deactivation}
When the monoester formed in the active site is re-acylated before it diffuses away (pseudo-processivity caused by mass-transfer limitation \cite{mitchell2023}), one productive encounter can add more than one acyl group. With $p_M$ the probability that a newly formed M is esterified again before release and $p_D$ the analogous probability for D, the event rates per unit of the common factor $r$ become
\begin{align}
\dot Z&=-Z, & \dot M&=Z(1-p_M)-S_MM,\nonumber\\
\dot D&=Zp_M(1-p_D)+S_MM(1-p_D)-S_DD, & \dot T&=Zp_Mp_D+S_MMp_D+S_DD,\nonumber\\
\Sigma&\equiv\dot c=Z(1+p_M+p_Mp_D)+S_MM(1+p_D)+S_DD,\label{eq:proc}
\end{align}
which reduces exactly to Eq.~(\ref{eq:fp}) for $p_M=p_D=0$. Enzyme deactivation inside a batch is described by first-order decay of the active enzyme, $r\to r\,e^{-k_dt}$. Both mechanisms therefore nest the original model: PPK-FIP can only differ from it where the data support $p_M$, $p_D$ or $k_d$ different from zero, which is what makes them testable.

\subsection{Exact Time Reparameterisation}
All reactions share the denominator of Eq.~(\ref{eq:dyn}); therefore the composition is a function of $c$ alone (Eq.~\ref{eq:proc} divided by $\Sigma$), and time follows by quadrature:
\begin{equation}
k_E\,\tau(c)=J_0(c)+F_AJ_1(c)+F_{AZ}J_Z(c)+F_{AM}J_M(c)+F_{AD}J_D(c),\qquad \tau=\frac{1-e^{-k_dt}}{k_d},\label{eq:tau}
\end{equation}
with $J_0=\int_0^c S/((A_0-c')\Sigma)\,dc'$, $J_1=\int_0^c dc'/\Sigma$ and $J_X=\int_0^c X/\Sigma\,dc'$ ($X=Z,M,D$). For $p_M=p_D=0$, $J_0=\ln[A_0/(A_0-c)]$. The four composition equations and five integrals are integrated once on a grid refined towards complete conversion (fourth-order Runge--Kutta compiled with Numba \cite{lam2015}); predictions at any time follow by monotone interpolation of Eq.~(\ref{eq:tau}). The rate constants $k_E$, $F_\ast$ and $k_d$ then enter only through a scalar product, so changing them costs well below a millisecond.

Figure~\ref{fig:bench} quantifies what this buys. For one evaluation of the most expensive model (Bornadel design, processive model with deactivation), adaptive ODE solvers need between 19\,ms (LSODA at a loose tolerance) and \SolverRadauMs\,ms (Radau at tolerance $10^{-9}$); the exact reparameterisation needs 11\,ms when the composition path must be recomputed (new $S_M$, $S_D$, $p_M$, $p_D$) and \SolverCachedMs\,ms when only $k_E$, $F_\ast$ or $k_d$ change, at a maximum error of $\SolverErr$ molar\,\% against a Radau reference at tolerance $10^{-12}$. The speed-up of one to three orders of magnitude at unchanged accuracy is not a convenience but a precondition of the study: the evaluation programme of Steps~1--6 (30 multistarts per fit, \NLotoFolds{} cross-validation folds times every method and ablation variant, and \NSimReps{} Monte-Carlo replicates per scenario) would otherwise be prohibitive.
\begin{figure}[!htbp]\centering\incfig[0.75\linewidth]{fig05_solver_benchmark}\caption{Cost and accuracy of one model evaluation (Bornadel design, processive model with deactivation): adaptive ODE solvers versus the exact reparameterisation, with the composition path recomputed (a new $S_M,S_D,p_M,p_D$) or cached (new $k_E$, $F_\ast$, $k_d$). Maximum error relative to Radau at tolerance $10^{-12}$.}\label{fig:bench}\end{figure}

\subsection{Fingerprint-Informed Prior, Model Averaging and Meta-Validation}
The two-step method is the limit of a Bayesian estimator in which $(\ln S_M,\ln S_D)$ have an infinitely tight prior at the fingerprint estimate, whereas a joint fit of all parameters corresponds to a flat prior. PPK-FIP uses the intermediate, calibrated choice
\begin{equation}
\ln S_X\sim\mathcal N\!\left(\ln\hat S_X^{\,\mathrm{fp}},\ (\lambda\,\delta_X/1.96)^2\right),\quad p_M,p_D\sim\mathcal N^+(0,s_p^2),\quad \log_{10}k_d\sim\mathcal N(-3,1.5^2),\label{eq:prior}
\end{equation}
where $\delta_X$ is the relative 95\,\% half-width of the fingerprint estimate (computed on the training data of each fold), $\lambda$ widens it, and $s_p$ shrinks processivity towards the original model. The MAP estimate is obtained by appending prior residuals $\hat\sigma(u-\mu)/s$ to the least-squares residual vector ($\hat\sigma$ = residual SD of the unpenalised fit) and minimising with a trust-region reflective algorithm from warm starts provided by the nested unpenalised family (joint; +$k_d$; +processivity; +both). Two variants are fitted, with the paper's simplified set of lumped constants and with the full set, and their predictions are averaged with BIC weights $w_i\propto e^{-\mathrm{BIC}_i/2}$ \cite{hoeting1999,burnham2002}. The hyper-parameters $(\lambda,s_p)\in\{1,3\}\times\{0.05,0.2\}$ are selected for each case by \emph{leave-one-case-out meta-validation}: the pair that minimises the mean cross-validated error ratio (PPK-FIP/original) over the \emph{other five} case studies is used, so no information from the target case enters the choice.

\subsection{Evaluation Protocols}
The protocols are listed in the order of the six steps of Table~\ref{tab:roadmap}.
\textbf{In-sample (paper protocol; Steps 1--2).} All methods are fitted to the full data of each case with 30 multistarts; we report $F_\mathrm{obj}$, BIC (paper convention), $\mathrm{BIC}_\mathrm{tot}$ and AICc \cite{hurvich1989}, and examine the residuals and the composition path against $c$.
\textbf{Leave-one-time-point-out cross-validation (LOTO; Step 3).} For each of the \NLotoFolds{} sampling times, all species measured at that time are removed and every method, including both steps of the original method, is re-estimated from scratch on the remaining data; the held-out values are then predicted. This is the primary out-of-sample criterion.
\textbf{Forward extrapolation (Step 3).} Each method is trained on the observations with $t\le0.5\,t_\mathrm{max}$ and predicts the rest. With a single split per case this protocol characterises behaviour beyond the data window rather than ranking methods statistically; it is also the protocol that separates mechanistic from data-driven models (Step 6).
\textbf{Ablation (Step 4).} One component of PPK-FIP at a time is removed (processivity, deactivation, the prior, the averaging, the meta-selection), and every variant is re-estimated in every LOTO fold. A multi-task extension in which profiles of the same system share the composition path (L1), additionally the lumped constants with hierarchical shrinkage (L1h) or everything (L2) is evaluated in the same way, with the level chosen by leave-one-case-out meta-validation.
\textbf{Cross-condition transfer (Step 5).} The composition path is learned from the other conditions of the same system (other temperature, water activity or batch) and the held-out profile is predicted against $c$, with and without processivity.
\textbf{Monte-Carlo truth recovery (Step 5).} Synthetic data are generated with the sampling design and residual noise of three real cases from three different truths (S1: the original model; S2: pseudo-processive; S3: deactivating), \NSimReps{} replicates each, and estimation error is measured against the noise-free truth on a dense time grid.
\textbf{Black-box baselines and hybrid (Step 6).} A Gaussian process (Matérn-5/2 + white noise) \cite{rasmussen2006}, a five-member MLP ensemble, the pretrained tabular foundation model TabICL (v2 regressor, in-context learning) \cite{qu2025tabicl} and a neural ODE trained through a Runge--Kutta solver \cite{chen2018} receive time, the condition descriptor (water activity or batch) and the known initial state. The hybrid PPK-FIP-GP adds to the mechanistic prediction a Gaussian-process discrepancy of Kennedy--O'Hagan type, fitted for each species to the training residuals, which reverts to zero away from the data; the same discrepancy is added to the original model for comparison.

\subsection{Statistics and Software}
Paired differences of squared prediction errors are tested with the Wilcoxon signed-rank test \cite{wilcoxon1945}, Holm-corrected across the six cases \cite{holm1979}; a more conservative test treats each fold rather than each observation as the unit. All computations used Python (NumPy, SciPy \cite{virtanen2020}, scikit-learn, PyTorch) on an NVIDIA DGX Spark (GB10) workstation; all numbers in this paper are generated from the result files by script.
